\section{Software implementation}
\label{sec:software}

\subsection{Package organisation and tensor interface}

The package is organised around a small public interface. \texttt{DiffusionConfig}
stores model and optimisation settings. \texttt{DDPM} owns the conditional U-Net,
forward diffusion process, and selected sampler. \texttt{StatisticsEvaluator}
computes ensemble diagnostics. \texttt{InverseProblem} estimates a scalar
conditioning parameter through a fitted DDIM generator. The package
does not contain a simulator for a particular physical system. Users create
trajectory--conditioning pairs with their own numerical solver or measurement
pipeline, then pass them to \texttt{DDPM.fit}.

The data interface has a batch axis, optional channel axis, and physical-time
axis. A scalar dataset has the shape \texttt{[batch, time]}. A multichannel
dataset has the shape \texttt{[batch, channels, time]}. Dataset members must be
ordered as \texttt{(data, conditioning)}. The number of generated channels and
conditioning channels is set in \texttt{DiffusionConfig}. This explicit
representation avoids implicit reshaping of physical variables during training.

\begin{pythonlisting}
from torch.utils.data import TensorDataset
from diffusionsde import DDPM, DiffusionConfig

config = DiffusionConfig(
    trajectory_length=256,
    n_output_channels=2,
    n_conditioning_channels=1,
)
ddpm = DDPM(config, sampler_type="ddim")
history = ddpm.fit(TensorDataset(data, conditioning))
\end{pythonlisting}

The method returns $M$ samples for each supplied conditioning input. If the input
contains $B$ conditions, the returned leading dimension is $BM$. For one generated
channel, the output shape is \texttt{[BM, N\_T]}; for $C_q>1$, it is
\texttt{[BM, C\_q, N\_T]}. The sampler flattens the first two axes. Downstream
code should reshape it only after retaining $B$ and $M$. The output is on the
model scale, so apply the inverse data transform before physical validation.

\subsection{Configuration and training}

Table~\ref{tab:config} gives the default configuration. The default values are
starting points, not physical constants. In particular, $N_K$ determines the
artificial noising process, whereas \texttt{trajectory\_length} is fixed by the
physical observation grid. Internal resolution is reduced by factors close to two.
The implementation accepts other trajectory lengths and interpolates decoder
features when needed. A length divisible by $2^{n_{\mathrm{layers}}-1}$ gives
exact resolution changes between U-Net levels. Dilation and the number of levels
should be selected from the correlation time and dominant frequencies of the data
rather than from $N_K$.

\begin{table}[t]
  \centering
  \caption{Selected \texttt{DiffusionConfig} defaults.}
  \label{tab:config}
  \small
  \begin{tabular}{@{}lll@{}}
    \toprule
    Group & Entry & Default \\
    \midrule
    Diffusion & number of diffusion steps $N_K$ & $1000$ \\
    Diffusion & schedule endpoints $\beta_1$, $\beta_{N_K}$ & $10^{-4}$, $0.02$ \\
    Data & physical trajectory length & $250$ \\
    Architecture & hidden width & $128$ \\
    Architecture & encoder levels & $4$ \\
    Architecture & kernel size & $3$ \\
    Architecture & generated and conditioning channels & $1$, $1$ \\
    Optimisation & learning rate & $10^{-4}$ \\
    Optimisation & batch size & $32$ \\
    Optimisation & epochs & $50$ \\
    Loss & form & MSE; Huber optional \\
    \bottomrule
  \end{tabular}
\end{table}

The trainer creates shuffled minibatches, applies the objective in
Eq.~\eqref{eq:noise_objective}, clips the global parameter-gradient norm at one,
and advances a cosine-annealed learning rate. A validation set is optional. When
it is present, the trainer records one held-out loss per epoch and restores the
weights from the epoch with the lowest validation loss. Training history is a
dictionary with \texttt{train\_loss}; it also contains \texttt{val\_loss},
\texttt{best\_val\_loss}, and \texttt{best\_epoch} when validation is used.

The checkpoint stores model weights, optimiser and scheduler states, loss
histories, validation information, and the configuration. The high-level
\texttt{DDPM.load} method restores the saved configuration and model weights for
generation. It does not restore an optimiser or scheduler. Sampler type and
sampler arguments are not stored in the checkpoint, so users must set them before
they create a \texttt{DDPM}. A workflow that resumes optimisation must use the
full trainer checkpoint state.

\begin{pythonlisting}
train_dataset = TensorDataset(x_train, c_train)
validation_dataset = TensorDataset(x_validation, c_validation)

history = ddpm.fit(
    train_dataset,
    val_dataset=validation_dataset,
)
ddpm.save("checkpoints/duffing.pt")

restored = DDPM(
    sampler_type="ddim",
    sampler_kwargs={"num_inference_steps": 200, "eta": 0.0},
)
restored.load("checkpoints/duffing.pt")
\end{pythonlisting}

\subsection{Network and sampler choices}

The conditional U-Net is described in Section~\ref{sec:conditional_diffusion}.
The implementation also provides \texttt{compute\_dilated\_rf}, which evaluates
the receptive field implied by a network depth, convolution width, and dilation
schedule. This check is useful before fitting long, correlated trajectories. For
causal applications, \texttt{causal=True} applies left-only convolution padding
and causal normalisation; a change to future input samples then cannot change an
earlier denoiser output.

\texttt{DDPMSampler} evaluates all $N_K$ reverse steps and draws posterior noise
at every nonterminal step. \texttt{DDIMSampler} can use fewer inference steps.
The \texttt{eta} parameter is the $\eta$ in Eqs.~\eqref{eq:ddim_sigma} and
\eqref{eq:ddim_update}. It is zero by default. With \texttt{eta=0}, a fixed
\texttt{initial\_noise} tensor makes a draw repeatable. Positive values add noise
during the reverse chain. The code below uses the preceding two-output-channel
configuration.

\begin{pythonlisting}
import torch

# c_query has shape [trajectory_length] because there is one condition channel.
initial_noise = torch.randn(
    32, config.n_output_channels, config.trajectory_length,
)

samples = ddpm.generate(
    c_query,
    n_samples=32,
    initial_noise=initial_noise,
    show_progress=False,
)
\end{pythonlisting}

The optional \texttt{require\_grad=True} flag retains automatic-differentiation
operations through DDIM generation. The \texttt{grad\_steps} argument can restrict
the graph to selected denoising steps, which reduces memory use at the cost of an
approximate gradient. It does not change the generated sample. A fixed initial
Gaussian tensor gives a repeatable inverse objective only when the DDIM sampler
uses $\eta=0$. This argument is specific to DDIM sampling and must not be sent to
the ancestral DDPM sampler.

\subsection{Validation diagnostics}

The evaluator compares two scalar ensembles of shape \texttt{[n\_samples, time]}.
It reports errors between the observed and generated mean and standard-deviation
curves. It also computes the per-trajectory autocorrelation function, an
ensemble-averaged Welch power spectrum \cite{welch1967psd}, and statistics of a
pooled tail marginal. The latter include the Wasserstein-1 distance
\cite{villani2009optimal}, skewness, and excess kurtosis. These diagnostics are
complementary. A generated ensemble can match pointwise means while having an
incorrect correlation time or power spectrum.

\begin{pythonlisting}
from diffusionsde import StatisticsEvaluator

evaluator = StatisticsEvaluator(reference, generated)
metrics = evaluator.compute_extended_metrics(
    dt=dt,
    max_lag=64,
    tail_fraction=0.5,
)
print(metrics["acf_mae"], metrics["wasserstein_1"])
\end{pythonlisting}

The \texttt{tail\_fraction} argument is important for relaxing systems. It limits
the autocorrelation, spectral, and marginal tests to the final part of each
trajectory, excluding an initial transient when stationarity is expected only at
long physical times. Autocorrelation is normalised separately for each trajectory
before ensemble averaging. The resulting metrics should be reported with their
definition and physical sampling interval, not only as a single aggregate score.

\subsection{Differentiable inverse estimation}

The inverse module treats a fitted generator as a conditional surrogate. It is
intended for one scalar parameter per optimisation run. The user supplies a
function that maps a proposed parameter value to a conditioning tensor. Let
$\mathcal{I}$ be the physical-time indices used by the solver, and let
$\overline{q}_n$ and $V_n$ denote the ensemble mean and sample variance at index
$n$. For an ensemble of $N_{\mathrm{ens}}$ trajectories, these are
\begin{align}
  \overline{q}_n
  &=
  \frac{1}{N_{\mathrm{ens}}}
  \sum_{s=1}^{N_{\mathrm{ens}}}q_n^{(s)},
  \label{eq:ensemble_mean}\\
  V_n
  &=
  \frac{1}{N_{\mathrm{ens}}-1}
  \sum_{s=1}^{N_{\mathrm{ens}}}
  \left(q_n^{(s)}-\overline{q}_n\right)^2 .
  \label{eq:ensemble_variance}
\end{align}
The solver generates an ensemble and minimises
\begin{equation}
  \mathcal{M}
  = w_{\mathrm{mean}}\frac{1}{\lvert\mathcal{I}\rvert}
    \sum_{n\in\mathcal{I}}
    \left(\overline{q}_n^{\mathrm{gen}}-\overline{q}_n^{\mathrm{obs}}\right)^2
  + w_{\mathrm{var}}\frac{1}{\lvert\mathcal{I}\rvert}
    \sum_{n\in\mathcal{I}}
    \left(V_n^{\mathrm{gen}}-V_n^{\mathrm{obs}}\right)^2 .
  \label{eq:moment_matching_loss}
\end{equation}
For \texttt{tail\_start=None}, $\mathcal{I}=\{0,\ldots,N_T-1\}$. Otherwise,
$\mathcal{I}$ is the selected final interval. The weights correspond to
\texttt{mean\_weight} and \texttt{variance\_weight}. The objective does not
compare autocorrelations or spectra. Those diagnostics remain necessary after an
inverse estimate has been obtained. Generated and observed trajectories must use
the same data scale.

The solver fixes a different initial Gaussian tensor for each restart and keeps it
fixed over that restart's optimisation iterations. With DDIM and $\eta=0$, this
makes the sampled realisation repeatable while preserving derivatives with respect
to the conditioning parameter. Two parameterisations are available: direct
physical values with a soft boundary penalty, and a sigmoid map that keeps the
estimate inside a specified interval.

\begin{pythonlisting}
from diffusionsde import InverseProblem, InverseProblemConfig

def build_conditioning(p, n_gen, n_time, device, run_idx):
    return p.reshape(1, 1, 1).expand(n_gen, 1, n_time)

ddpm.freeze()
config = InverseProblemConfig(n_iterations=100, n_gen=64)
problem = InverseProblem(
    ddpm, observed, build_conditioning, (p_min, p_max), config,
)
result = problem.solve(n_runs=8)
\end{pythonlisting}

Use a DDIM sampler with $\eta=0$ for this workflow; the inverse module does not
enforce this setting. It is a practical way to estimate a parameter from ensemble
moments, but it does not prove parameter identifiability. The conditioning design,
training coverage, choice of moments, and agreement of temporal diagnostics must
be assessed for each physical problem.
